\documentclass[11pt,a4paper]{article}

\usepackage[margin=0.95in]{geometry}
\usepackage{amsmath,amssymb}
\usepackage{booktabs}
\usepackage{array}
\usepackage{float}
\usepackage{hyperref}
\usepackage{enumitem}
\usepackage{xcolor}
\usepackage{microtype}

\hypersetup{
    colorlinks=true,
    linkcolor=blue!60!black,
    urlcolor=blue!60!black,
    citecolor=blue!60!black
}

\title{\textbf{FOML Project Report}\\[0.3em]
\large Structured Prediction for Dependency Parsing}

\author{
    \textbf{Murthy Manideep} (CS24BTECH11039) \and
    \textbf{Adinesh Reddy} (CS24BTECH11026)
}

\date{}

\begin{document}

\maketitle

\section{Introduction}

In this project, we built a dependency parser for the FoML 2026 Hackathon. Given a sentence of words $w_1, w_2, \dots, w_n$ (plus a dummy root $w_0 = \texttt{<ROOT>}$), we need to build a parse tree over the sentence. For each word $i \in \{1, \dots, n\}$, we predict:
\begin{itemize}[noitemsep]
    \item its \textbf{head} $h_i \in \{0, 1, \dots, n\} \setminus \{i\}$, where $0$ means the root of the sentence, and
    \item its \textbf{label} $\ell_i$, which is the grammar relation to its head (like \texttt{nsubj}, \texttt{obj}, \texttt{det}, \texttt{obl}, \texttt{nmod}, or \texttt{punct}).
\end{itemize}
Each prediction is saved as \texttt{HEAD|LABEL}. Accuracy is measured using \textbf{Labeled Attachment Score (LAS)}, where a word is right only if both its head and label are right. While testing our models, we also checked \textbf{Unlabeled Attachment Score (UAS)} (head accuracy alone) so we could tell whether mistakes came from wrong heads or wrong labels.

We wrote the entire codebase in Python using only \texttt{numpy} and Python's built-in modules, using averaged perceptrons, Chu-Liu/Edmonds tree building, graph-based parsing, and transition-based parsing.

\section{Data and Validation Setup}

The dataset has \textbf{12,439 training sentences} in \texttt{train.csv} and \textbf{2,077 test sentences} in \texttt{test.csv}. Only the words are given, with no POS tags or lemmas.

\subsection{Block Validation Split (\texttt{--check})}
In our first tries, we kept every 10th single sentence for validation. That gave higher validation scores than Kaggle because nearby sentences in the file often come from the same article and share the same names and words (only 3.9\% new words in validation vs.\ 6.3\% in \texttt{test.csv}).

To make validation closer to the real test set, our \texttt{--check} mode groups \texttt{train.csv} into \textbf{blocks of 100 neighboring sentences} and holds out every 10th block (90\% train, 10\% validation). That raised the unseen word rate to 5.0\%. We used this split to test our features, parser combinations, and rules.

\subsection{5-Fold Cross-Validation Inside the Pipeline}
Our later models (transition parsers and the tree labeler) use predictions from earlier models (taggers and the MST parser) as input features. If a tagger predicted on the same sentences it trained on, those tags would be almost 100\% right on train but noisier on test, causing the parsers to overfit.

To avoid this, every earlier stage runs \textbf{5-fold cross-validation} on the training set (and builds a separate word stats table inside each fold). That way, the parsers train on realistic tag and MST mistakes that look like what they will see on test data.

\section{Pipeline and Models}

Our final system has six steps across six Python files: \texttt{perceptron.py}, \texttt{data\_utils.py}, \texttt{tagger.py}, \texttt{parser.py}, \texttt{combiner.py}, and \texttt{train\_and\_predict.py}.

\subsection{Averaged Perceptron}
All of our taggers, parsers, and the tree labeler use an \textbf{averaged perceptron with a small margin}. For an input $x$ and a choice $y$, the score is:
\[
s(x, y) = \mathbf{w}^\top \phi(x, y).
\]
Feature strings are hashed into a fixed-size weight table ($2^{20}$ rows for taggers and labeler, $2^{21}$ rows for transition parsers, and $2^{22}$ rows for Graph-MST). During training, we add a small penalty $\Delta(y, y^*) > 0$ to wrong choices ($\Delta = 0.5$ for taggers and labeler, $\Delta = 1.0$ for parsers) so the correct answer stays clearly above close wrong choices:
\[
\hat{y} = \arg\max_{y} \big(\mathbf{w}^\top \phi(x, y) + \Delta(y, y^*)\big).
\]
Whenever $\hat{y} \neq y^*$, we update the weights by 1:
\[
\mathbf{w} \leftarrow \mathbf{w} + \phi(x, y^*) - \phi(x, \hat{y}).
\]
At the end of training, we average the weights across all training steps, $\bar{\mathbf{w}} = \frac{1}{T}\sum_{t=1}^{T}\mathbf{w}^{(t)}$. Averaging keeps the weights smooth and helps a lot on unseen test sentences.

\subsection{Step 1: Word Cleaning, Stats, and Learned Word Groups}
In \texttt{data\_utils.py}, we clean the words and collect simple word features before training:
\begin{itemize}[noitemsep]
    \item \textbf{Word cleaning:} We lowercase words, remove apostrophes (because the data has both ``we ve'' and ``n't''), and replace URLs, emails, 4-digit years, percentages, money, times, and numbers with common tokens (\texttt{<URL>}, \texttt{<YEAR>}, \texttt{<NUMBER>}, etc.).
    \item \textbf{Stems and word patterns:} We take the first 5 letters as a simple stem, and build character patterns using capitals (\texttt{ALL\_CAP}, \texttt{INIT\_CAP}, \texttt{MID\_CAP}, \texttt{MIX\_CAP}, \texttt{LOW}), digits, hyphens, and the last 3 letters.
    \item \textbf{Word stats table:} From the training fold only, we save each word's most common tag, its top-2 common tags, its coarse tag group, and its most common left and right neighbor groups. For rare or unseen words, we fall back to the 5-letter stem and then to the character pattern.
    \item \textbf{Learned word groups:} We learn three small word sets directly from \texttt{train.csv}:
    \begin{enumerate}[noitemsep]
        \item \emph{Punctuation forms:} tokens labeled \texttt{punct} at least 98\% of the time (seen $\ge 3$ times). This avoids marking symbols like \texttt{\&} (\texttt{cc}), \texttt{\$} (money head), \texttt{\%} (noun), and \texttt{:)} (\texttt{discourse}) as \texttt{punct}.
        \item \emph{Standalone words:} common noun/pronoun-like words ($\ge 60\%$ nominal labels) that almost never have determiner, adjective, or compound children ($\le 2\%$), which finds pronouns and standalone names.
        \item \emph{Possessive words:} \texttt{case} words whose head is \texttt{nmod:poss} at least 80\% of the time.
    \end{enumerate}
\end{itemize}

\subsection{Steps 2--4: Five Sequence Taggers (3 Seeds Averaged)}
Because POS tags are not given in the data, we train five taggers in \texttt{tagger.py} to give useful grammar hints to the parsers. Each tagger uses 5-fold CV on train and \textbf{3-seed weight averaging} on test:
\begin{enumerate}[noitemsep]
    \item \textbf{Word-Class Tagger (UPOS, 6 epochs):} Builds 18 word classes (\texttt{NOUN}, \texttt{PROPN}, \texttt{PRON}, \texttt{VERB}, \texttt{VERB\_VBN}, \texttt{VERB\_INF}, \texttt{ADJ}, \texttt{ADV}, \texttt{ADP}, \texttt{DET}, \texttt{AUX}, \texttt{COP}, \texttt{CCONJ}, \texttt{SCONJ}, \texttt{NUM}, \texttt{PART}, \texttt{INTJ}, \texttt{PUNCT}) from the training trees and learns to predict them on new sentences.
    \item \textbf{Left-to-Right Tagger (6 epochs):} Reads left to right and guesses each word's label and head direction (like \texttt{det>}, \texttt{obj<}, \texttt{root}).
    \item \textbf{Right-to-Left Tagger (6 epochs):} Reads right to left so it can see words coming after the current word.
    \item \textbf{Combined Tagger (6 epochs):} Takes the left-to-right tags, right-to-left tags, whether they agree, UPOS tags, and distances to the nearest root, verb, and noun.
    \item \textbf{Head-Distance Tagger (5 epochs):} Predicts how far away each word's head is and what type of word the head is (for example \texttt{+1:OBJ}, \texttt{-2:CLAUSE}, \texttt{ROOT:ROOT}).
\end{enumerate}

\subsection{Step 5: Graph-MST Parser (4 Seeds Averaged)}
In \texttt{parser.py}, \texttt{GraphMSTParser} scores every possible edge $(h \to m)$ in a sentence at once using \textbf{42 fast numpy feature templates} (head and child words, stems, tags, UPOS, word stats, character patterns, direction and distance, counts of verbs/commas/conjunctions in between, neighbor tags, and head-distance tag matches).

Given the $(n+1)\times(n+1)$ edge score matrix $S$, we pick the best single-root tree rooted at token $0$:
\[
T^* = \arg\max_{T \in \mathcal{T}_1(n)} \sum_{m=1}^{n} S[\text{head}_T(m), m]
\]
using the recursive \textbf{Chu-Liu/Edmonds (CLE)} algorithm, making sure only one word attaches to $0$ (single root) and punctuation tokens cannot be heads of normal words. On train we run 5-fold CV (8 epochs), and on test we average the score matrices from \textbf{4 random seeds} before picking the tree. The MST output helps twice: it gives top-2 head hints and confidence levels (\texttt{HIGH}, \texttt{MID}, \texttt{LOW}) to every transition parser, and it also votes in the final tree selection.

\subsection{Step 6: 18 Transition Parsers, 4-Seed Tree Labeler, and Voting}

\paragraph{18 Transition Parsers (3 Types).}
Transition parsers build the tree one move at a time using a stack and a buffer. Because parts of the tree are already built as they go, they can check already-attached children, grandchildren, child counts, prepositions, auxiliaries, and child labels. To get better accuracy and lower variance, we train \textbf{18 parsers split evenly across three types}:
\begin{itemize}[noitemsep]
    \item \textbf{Arc-Hybrid} (6 parsers: 4 left-to-right, 2 right-to-left): makes right edges bottom-up and left edges between stack and buffer.
    \item \textbf{Arc-Eager} (6 parsers: 4 left-to-right, 2 right-to-left): makes right edges top-down with a \texttt{REDUCE} move and uses already-assigned parent labels.
    \item \textbf{Arc-Standard} (6 parsers: 4 left-to-right, 2 right-to-left): makes both left and right edges between the top two stack words while looking at the next buffer word.
\end{itemize}
Each parser trains for 12 epochs using \textbf{mistake-recovery training}: instead of always forcing the perfect training move, we let the parser follow its own chosen move (50\% of the time in epoch 0 and 90\% in epochs 1--11) and teach it whichever legal move loses the fewest remaining true edges from that state. This helps the parser keep going cleanly even after an early mistake. Right-to-left parsers run on reversed sentences with flipped tags and MST hints, then flip their predicted tree back.

\paragraph{4-Seed Full-Tree Labeler.}
Once the full tree is built, \texttt{FullTreeLabeler} scores relation labels for every edge $(h \to m)$ using full-sentence tree context: grandparent tags and UPOS, left and right siblings, whether head $h$ already has another subject or object child, preposition/mark/auxiliary/conjunction children, head word stats, relative-clause hints (\texttt{acl:relcl} vs.\ \texttt{acl}/\texttt{advcl}/\texttt{ccomp}), nearby punctuation, and child counts. Each of the 4 seeds trains for 6 epochs on a \textbf{70\% true trees / 30\% predicted MST trees mix} so it still predicts good labels when predicted heads have small mistakes.

\paragraph{Weighted Voting and 4 Cleanup Rules.}
All 18 transition parsers vote on edges using their parser weight and score margin, along with votes from the 4-seed Graph-MST parser (top-3 heads plus margin bonus). Running Chu-Liu/Edmonds on the combined score matrix guarantees one valid single-root tree per sentence. Labels are then chosen by combining parser label votes (giving higher weight when the parser also picked the same head) and scores from the 4-seed \texttt{FullTreeLabeler}. Finally, we run four simple rules:
\begin{enumerate}[noitemsep]
    \item \texttt{punct}: learned punctuation words always get the \texttt{punct} label.
    \item \texttt{flat}: keeps multi-word capitalized proper names connected with \texttt{flat}.
    \item \texttt{passive}: changes \texttt{nsubj} to \texttt{nsubj:pass} when the head verb has an \texttt{aux:pass} child (and no \texttt{cop} or \texttt{obj}).
    \item \texttt{obl\_nmod}: fixes \texttt{obl} vs.\ \texttt{nmod} and possessive \texttt{nmod:poss} using learned possessive words and whether the head is a noun or verb.
\end{enumerate}
We also wrote an optional logistic regression combiner in \texttt{combiner.py} (using Newton's method in \texttt{numpy}) for \texttt{--check} mode, while our final submission uses the direct weighted vote and cleanup rules.

\section{Experiments and Results}

\subsection{5-Fold CV Accuracy of Tagger Steps}
Table~\ref{tab:cv} shows the 5-fold CV accuracy of our taggers on all 12,439 training sentences. Combining the left-to-right and right-to-left taggers in Step~3 improves accuracy from 84.98\% / 84.78\% to 85.34\%.

\begin{table}[H]
\centering
\caption{5-fold CV accuracies on all 12,439 training sentences.}
\label{tab:cv}
\begin{tabular}{@{}lc@{}}
\toprule
\textbf{Pipeline Step} & \textbf{5-Fold CV Accuracy} \\
\midrule
Step 1: Word-Class Tagger (18 UPOS classes) & 94.79\% \\
Step 2a: Left-to-Right Tagger & 84.98\% \\
Step 2b: Right-to-Left Tagger & 84.78\% \\
Step 3: Combined Both-Direction Tagger & 85.34\% \\
Step 4: Head-Distance Tagger & 80.34\% \\
\bottomrule
\end{tabular}
\end{table}

\subsection{Progress Across Iterations on Kaggle}
Table~\ref{tab:kaggle} shows how our Kaggle Public Leaderboard score (LAS) improved across iterations.

\begin{table}[H]
\centering
\caption{Score progress across iterations on the Kaggle Public Leaderboard (LAS).}
\label{tab:kaggle}
\begin{tabular}{@{}cl>{\raggedright\arraybackslash}p{7.6cm}c@{}}
\toprule
\textbf{Iter.} & \textbf{Model} & \textbf{What We Changed} & \textbf{Public LAS} \\
\midrule
1 & Baseline & 5 Arc-Hybrid parsers + simple vote & 0.7803 \\
2 & + Second Pass & Added second-pass tagger, 7 parsers & 0.7811 \\
3 & + Stats \& Tags & Both-direction taggers, word stats, head-distance tagger & 0.8038 \\
4 & + Graph-MST & Added Graph-MST parser to guide transition parsers & 0.8119 \\
5 & + 3 Parser Types & Added UPOS, 12 parsers (Hybrid/Eager/Standard), full-tree labeler, CLE vote & 0.8199 \\
6 & + Learned Groups & Learned punct/standalone/possessive sets from train, apostrophe fix, 4 rules & 0.8255 \\
7 & \emph{Experiment} & Variable step updates, beam-4 tagger, extra 2nd-order step, cluster templates \emph{(reverted)} & 0.8235 \\
8 & + Multi-Seed & Reverted Iter.\ 7; 3-seed taggers/MST/labeler, 15 balanced parsers, extra labeler features & 0.8288 \\
\textbf{9} & \textbf{Final Model} & \textbf{4-seed MST, 4-seed labeler, 18 parsers (6+6+6), balanced MST/labeler vote weights, extra labeler features} & \textbf{Final} \\
\bottomrule
\end{tabular}
\end{table}

\subsection{Things We Tried That Did Not Work}
In Iteration~7, our score went down from \texttt{0.8255} to \texttt{0.8235}. Checking why showed four problems:
\begin{enumerate}[noitemsep]
    \item \textbf{Variable step-size updates:} Dividing perceptron updates by the number of active features made updates too small on rare words and shrank MST margins, breaking our \texttt{HIGH}/\texttt{MID}/\texttt{LOW} confidence thresholds. Going back to simple step-size 1 updates fixed this immediately.
    \item \textbf{Beam-4 early stopping in taggers:} Stopping a sentence early when the true sequence dropped out of the top-4 beam skipped learning on the second half of long sentences.
    \item \textbf{Extra second-order search step:} Running a separate greedy second-order scorer on top of the parser votes added noise and sometimes broke good agreed edges.
    \item \textbf{Extra cluster features:} Adding large word-cluster templates into fixed-size hash tables caused more hash overlaps without helping accuracy.
\end{enumerate}
Removing those changes and using \textbf{multi-seed averaging}, \textbf{balanced 3-type parser ensembles}, and \textbf{better full-tree labeler features} raised our score from \texttt{0.8255} to \texttt{0.8288} and our final selected model.

\section{How to Run}

\begin{enumerate}[noitemsep]
    \item Put \texttt{train.csv}, \texttt{test.csv}, and \texttt{sample.csv} inside \texttt{data/}.
    \item Install \texttt{numpy}: \texttt{pip install numpy}.
    \item Run \texttt{python3 -u train\_and\_predict.py}.
\end{enumerate}
Final predictions are written to \texttt{outputs/submission.csv} with fixed seeds and \texttt{PYTHONHASHSEED=0}.

\section{Team Member Contributions}
\begin{table}[H]
\centering
\begin{tabular}{@{}p{4.0cm}>{\raggedright\arraybackslash}p{11.0cm}@{}}
\toprule
\textbf{Team Member} & \textbf{Work Done} \\
\midrule
\textbf{Murthy Manideep} \newline (\texttt{CS24BTECH11039}) &
Data reading and word cleaning, word stats table and learned word sets, averaged perceptron, all 5 sequence taggers (UPOS, left-to-right, right-to-left, combined, head-distance), Graph-MST parser with Chu-Liu/Edmonds tree builder, error analysis, and model training runs. \\[0.6em]
\textbf{Adinesh Reddy} \newline (\texttt{CS24BTECH11026}) &
All 3 transition parser types (Arc-Hybrid, Arc-Eager, Arc-Standard) with mistake-recovery training, right-to-left sentence flipping, 4-seed full-tree labeler and features, ensemble voting, 4 post-processing rules, 90/10 block validation checks, and report/README writing. \\
\bottomrule
\end{tabular}
\end{table}

\section{Conclusion}

Overall, simple linear structured models with careful feature engineering worked very well for dependency parsing. The biggest gains came from avoiding overfitting with 5-fold CV inside the pipeline, combining a global Chu-Liu/Edmonds MST parser with 18 transition parsers across three systems and both directions, training the full-tree labeler on a mix of true and predicted trees, and learning punctuation and possessive word sets directly from \texttt{train.csv}.

\end{document}
